\par\Needspace{9\baselineskip}\bigskip
\pdfbookmark[3]{Exposition, [23]}{pallas-II-article23}
\begin{center}[23.]\end{center}
\begingroup
\setlength{\lineskip}{4pt}\setlength{\lineskiplimit}{4pt}
\noindent\hspace*{\parindent}This method gives the elements \(\AscendingNode+\delta\AscendingNode,\ i+\delta i,\ \varphi+\delta\varphi,\ \tilde\omega+\delta\tilde\omega,\ N+\delta N,\ n+\delta n\). In the expansion of the instantaneous variation of \(n+\delta n\), the coefficients of all the terms for which \(i=0\) become zero. It follows that neither the mean motion nor the major axis contains a secular variation in the strict sense, at least if we stop at the first power of the perturbing force. These elements will also contain no periodic terms independent of \(L\). To prove this directly, let us take up the formula \(m\mathrm{d}\dfrac{1}{a}\bigl(\text{or rather }m\mathrm{d}\dfrac{1}{a+\delta a}\bigr)=2(\xi\mathrm{d}x+\eta\mathrm{d}y+\zeta\mathrm{d}z)\) again. It is easy to see that if we set
\[
m\left\{\frac{xx'+yy'+zz'}{r'^{8}}-\frac{1}{\rho}\right\}=\Omega
\]
\SourceLeaf{473}the forces \(\xi,\eta,\zeta\) are expressed by the partial differentials
\[
\xi=\mu\left(\frac{\mathrm{d}\Omega}{\mathrm{d}x}\right)\mathpunct{\raisebox{.45ex}{,}}\qquad
\eta=\mu\left(\frac{\mathrm{d}\Omega}{\mathrm{d}y}\right)\mathpunct{\raisebox{.45ex}{,}}\qquad
\zeta=\mu\left(\frac{\mathrm{d}\Omega}{\mathrm{d}z}\right)\mathpunct{\cdot}
\]
Thus \(\mathrm{d}\dfrac{2\mu}{a+\delta a}\) will be the differential of \(m\Omega\) when the quantities relating to the perturbing planet are treated as constant. It is now clear that \(\Omega\) is a \textit{periodic} function of \(\boldsymbol{L}\) and \(\boldsymbol{L}'\). If we therefore expand it in terms of the form \(k\cos(iL+i'L')\) and \(k\sin(iL+i'L')\), its differential, with \(L'\) treated as constant, will contain no term independent of \(L\). And since we have
\[
\begin{aligned}
n+\delta n&=\BareRoot{m}.(a+\delta a)^{-\frac{3}{2}}\\
\mathrm{d}(n+\delta n)&=-\tfrac{3}{2}\BareRoot{m}.(a+\delta a)^{-\frac{5}{2}}\mathrm{d}(a+\delta a)=\tfrac{3}{2}\BareRoot{\dfrac{m}{a+\delta a}}\cdot\mathrm{d}\frac{1}{a+\delta a}\mathpunct{\raisebox{.45ex}{,}}
\end{aligned}
\]
we obtain, considering only the first-order terms,
\[
\mathrm{d}(n+\delta n)=3\frac{\mu}{\BareRoot{ma}}\cdot\mathrm{d}\Omega,
\]
with the differential \(\mathrm{d}\Omega\) taken with respect to \(L\). It follows that the variation \(\mathrm{d}(n+\delta n)\) has the form indicated above.
\par
If \(\bar n\) is the constant introduced by integrating \(\mathrm{d}(n+\delta n)\), it is clear that what corresponds to the mean motion, that is, the integral \(\int n\mathrm{d}t\), will consist
\begin{enumerate}\setlength{\itemsep}{0pt}\setlength{\parsep}{0pt}
\item[1)] of periodic terms
\item[2)] of a term proportional to the\raisebox{.45ex}{-} time, \(=\bar n t\)
\item[3)] of a constant part.
\end{enumerate}
\par
If we therefore combine this mean motion with the variable epoch, which will itself also consist
\begin{enumerate}\setlength{\itemsep}{0pt}\setlength{\parsep}{0pt}
\item[1)] of periodic terms
\item[2)] of a secular variation that we shall denote by \(\bar\varepsilon t\)
\item[3)] of a constant part,
\end{enumerate}
\noindent the mean longitude will take the form \(\bar L+(\bar n+\bar\varepsilon)t+P\), where \(P\) denotes the periodic parts.
\par
This leaves six constant terms, which we shall call the \textit{mean values of the elements}. Combining them with the \SourceLeaf{474}periodic parts and the secular variations gives the complete values of the elements at any instant.
\par
For the constant parts of the elements \(\AscendingNode,\ i,\ \varphi,\ \tilde\omega\), from which the perturbations are calculated, it is appropriate to choose the mean values themselves, that is, the constants introduced by integration, and for [those of] \(L\) and \(n\) the quantities \(\bar L\) and \(\bar n+\bar\varepsilon\). It is clear, however, that in practice this will require two calculations. After [making] the first calculation using any approximate elements, we shall determine by known methods the values to assign to the six constants to obtain the best agreement with the observations, then repeat the calculation with these values. Unless the first values were too far from the true ones, we shall certainly have no need to make the calculation for the third time.
\par\endgroup
